% =============================================================================
%  Apéndice A -- Instrumento de la encuesta
%  -------------------------------------------------------------------------
%  NOTA PARA AXEL:
%  - Este es un BORRADOR COMPLETO del instrumento, listo para discutir con tu
%    asesor antes de aplicarlo. Tres decisiones de diseño que debes validar:
%    (1) escala Likert de 5 puntos con etiqueta en cada punto;
%    (2) los ítems L7-L9 usan viñetas (mini-escenarios) en vez de preguntas
%        directas, para reducir deseabilidad social;
%    (3) las abiertas A1-A4 están ordenadas de lo concreto a lo abstracto.
%  - Marcas [DECIDIR]: puntos donde la decisión es tuya o de tu asesor.
%  - Antes de aplicar: pilotear con 5-8 personas y medir tiempo de respuesta
%    (objetivo: menos de 12 minutos).
% =============================================================================
\chapter{Instrumento de la encuesta}
\label{ap:encuesta}

\section*{Ficha técnica del instrumento}

\begin{description}
  \item[Título.] Percepciones sobre lenguaje, tecnología y trato diferenciado.
  \item[Población objetivo.] Personas de 18 a 30 años residentes en México.
    [DECIDIR: confirmar rango de edad con el asesor.]
  \item[Modo de aplicación.] Formulario en línea, autoadministrado, anónimo.
  \item[Duración estimada.] 10--12 minutos.
  \item[Estructura.] Cuatro secciones: consentimiento informado, datos
    sociodemográficos, percepciones (escala Likert) y preguntas abiertas.
\end{description}

\section*{Sección 1. Consentimiento informado}

\begin{quote}
Esta encuesta forma parte de una investigación de tesis de licenciatura en
Ciencias de la Computación (UNAM) sobre la relación entre el estilo de
escritura y el trato que las personas reciben de sistemas automatizados
(por ejemplo, filtros de currículums o moderación de contenido en plataformas).

Tu participación es \textbf{voluntaria y anónima}: no se solicita tu nombre ni
ningún dato que permita identificarte. Las respuestas se utilizarán únicamente
con fines académicos y se reportarán de forma agregada; las respuestas a
preguntas abiertas podrán citarse textualmente en la tesis sin atribución.
Puedes abandonar el cuestionario en cualquier momento sin consecuencia alguna.

Responsable: Axel Casas Espinosa (correo de contacto: [DECIDIR]).
Asesor: M.\ en C.I.C.\ Manuel Soto Romero, Facultad de Ciencias, UNAM.

$\square$ He leído lo anterior y acepto participar.
\end{quote}

\section*{Sección 2. Datos sociodemográficos}

\begin{enumerate}[label=\textbf{D\arabic*.}]
  \item Edad (años cumplidos): \rule{2.5cm}{0.4pt}
  \item Estado de residencia: \rule{5cm}{0.4pt}
  \item ¿Tu localidad es principalmente...? \quad
    $\square$ Urbana \quad $\square$ Suburbana \quad $\square$ Rural
  \item Último nivel educativo concluido: \quad
    $\square$ Secundaria o menos \quad $\square$ Bachillerato \quad
    $\square$ Licenciatura \quad $\square$ Posgrado
  \item ¿Actualmente estudias, trabajas, ambas o ninguna? \quad
    $\square$ Estudio \quad $\square$ Trabajo \quad $\square$ Ambas \quad
    $\square$ Ninguna
  \item Máximo nivel educativo concluido por tu madre o tutor(a) principal:
    \quad $\square$ Primaria o menos \quad $\square$ Secundaria \quad
    $\square$ Bachillerato \quad $\square$ Licenciatura o más \quad
    $\square$ No sé
  \item Pensando en el hogar donde creciste, ¿con cuál de estas frases lo
    describirías mejor? \quad
    $\square$ El dinero no alcanzaba ni para lo básico \quad
    $\square$ Alcanzaba justo para lo básico \quad
    $\square$ Alcanzaba para lo básico y algo más \quad
    $\square$ Vivíamos con holgura
\end{enumerate}

\noindent\textit{Nota metodológica: D6 y D7 son aproximaciones (proxies) de
nivel socioeconómico de origen, menos invasivas que preguntar ingreso. Si el
asesor prefiere un instrumento estandarizado, considerar la regla AMAI
[DECIDIR].}

\section*{Sección 3. Percepciones (escala Likert)}

\noindent Indica qué tan de acuerdo estás con cada afirmación: \\
\textbf{1} = Totalmente en desacuerdo \quad \textbf{2} = En desacuerdo \quad
\textbf{3} = Ni de acuerdo ni en desacuerdo \quad \textbf{4} = De acuerdo
\quad \textbf{5} = Totalmente de acuerdo

\subsection*{Bloque I. Discriminación lingüística en general}

\begin{enumerate}[label=\textbf{L\arabic*.}]
  \item En México, la forma de hablar o escribir de una persona influye en
    cómo la tratan los demás.
  \item He visto que se juzgue a alguien como menos inteligente o menos capaz
    por su manera de expresarse.
  \item Yo he modificado mi forma de hablar o escribir para causar una mejor
    impresión en contextos formales (escuela, trabajo, trámites).
  \item Alguna vez me he sentido juzgado(a) o tratado(a) diferente por mi
    forma de hablar o escribir.
\end{enumerate}

\subsection*{Bloque II. Tecnología y trato diferenciado}

\begin{enumerate}[label=\textbf{L\arabic*.},resume]
  \item Creo que los sistemas automatizados (filtros de currículums,
    moderación en redes, atención por chat) tratan igual a todas las personas
    sin importar cómo escriban.\footnote{Ítem invertido: se redacta en sentido
    contrario a los demás para detectar respuesta automática. Recodificar al
    analizar.}
  \item Me preocupa que una computadora evalúe mis textos (por ejemplo, mi
    currículum) sin intervención humana.
\end{enumerate}

\subsection*{Bloque III. Escenarios}

\noindent\textit{Lee cada situación y responde con la misma escala.}

\begin{enumerate}[label=\textbf{L\arabic*.},resume]
  \item Dos personas con la misma experiencia laboral envían su currículum
    a una empresa que usa un filtro automatizado. Una escribe de manera muy
    formal; la otra, de manera sencilla y directa, con algunos modismos.
    \textit{Creo que el filtro las evaluaría igual.}\footnotemark[\value{footnote}]
  \item Si un sistema automatizado rechazara mi solicitud (de empleo, beca o
    crédito), sospecharía que mi forma de escribir influyó en el resultado.
  \item Considero justo que una empresa use la ``calidad de redacción'' como
    criterio automático para filtrar candidatos, aunque el puesto no requiera
    escribir.
\end{enumerate}

\subsection*{Bloque IV. Experiencia directa}

\begin{enumerate}[label=\textbf{L\arabic*.},resume]
  \item He interactuado con sistemas automatizados que evalúan o filtran lo
    que escribo (postulaciones en línea, chatbots de soporte, moderación en
    redes sociales).
  \item En esas interacciones, he sentido que el resultado habría sido
    distinto si hubiera escrito de otra manera.
  \item Confío en que las decisiones automatizadas sobre textos pueden
    apelarse o revisarse por una persona si lo solicito.
\end{enumerate}

\section*{Sección 4. Preguntas abiertas}

\noindent\textit{Responde con tus propias palabras; no hay respuestas
correctas ni incorrectas. La extensión es libre.}

\begin{enumerate}[label=\textbf{A\arabic*.}]
  \item Cuéntanos alguna ocasión en la que sientas que a ti o a alguien
    cercano lo trataron diferente (mejor o peor) por su forma de hablar o
    escribir. ¿Qué pasó?
  \item ¿Has cambiado alguna vez tu forma de escribir al llenar una solicitud,
    postularte a un empleo o hablar con un chatbot o sistema automatizado?
    ¿Qué cambiaste y por qué?
  \item ¿Qué palabras, expresiones o formas de escribir crees que una empresa
    o un sistema automatizado vería ``mal'' en un currículum? Da ejemplos si
    puedes.
  \item Si pudieras decirle algo a quienes diseñan estos sistemas
    automatizados, ¿qué les dirías?
\end{enumerate}

\section*{Cierre}

\begin{quote}
Gracias por participar. Si conoces a otras personas de entre 18 y 30 años que
puedan responder esta encuesta, te agradecemos compartirla
[muestreo por bola de nieve].
\end{quote}

% -----------------------------------------------------------------------------
% Mapa instrumento -> análisis (no se imprime en la versión aplicada;
% conservar aquí como documentación metodológica para el Cap. 3):
%
%   D1-D7   -> caracterización de la muestra (Cap 4, tabla descriptiva)
%   D6-D7   -> variable de cruce socioeconómico para los bloques Likert
%   L1-L4   -> percepción de discriminación lingüística (humana)
%   L5-L9   -> percepción de discriminación tecnológica  (algorítmica)
%   L10-L12 -> exposición y confianza
%   L5, L7  -> ítems invertidos (control de aquiescencia; recodificar)
%   A1-A2   -> corpus: experiencias narradas (temas + sentimiento)
%   A3      -> corpus: léxico estigmatizado percibido -> INSUMO DIRECTO para
%              la tipología de variaciones del benchmark (Fase 3)
%   A4      -> corpus: expectativas normativas (Cap 6, discusión)
% -----------------------------------------------------------------------------
